% --- Archon physics formalization source begin ---
% archon:physics
% archon:covers PhyXMiniProblems/problem_phyx_mini_0057.lean
% archon:source-report reports/phyx_mini/problem_phyx_mini_0057.source.json

\chapter{Physics problem phyx\_mini\_0057}
\label{ch:PhyXMiniProblems_problem_phyx_mini_0057}

\paragraph{Problem source.}
\#\# Physical scenario

The objective lens of a microscope uses a planoconvex glass lens with the flat side facing the specimen. A real image is formed 160 mm behind the lens when the lens is 8.0 mm from the specimen.

\#\# Question

What is the radius of the lens’s curved surface?

\#\# Answer choices

- A: A: 1.4mm

- B: B: 8.6mm

- C: C: 3.8mm

- D: D: 7.4mm

\#\# Figure caption (auxiliary; use the image as primary evidence)

The image is a diagram illustrating the refraction of light through a lens. Here's a detailed description of the elements:

1. **Arrow Labelled "s" and "s'"**: 
   - An arrow at the bottom indicates an object distance \textbackslash{}( s = 8.0 \textbackslash{}, \textbackslash{}text\{mm\} \textbackslash{}).
   - An arrow at the top indicates an image distance \textbackslash{}( s' = 160 \textbackslash{}, \textbackslash{}text\{mm\} \textbackslash{}).

2. **Lens Properties**:
   - There is a translucent, lens-like shape denoting a refractive medium.
   - The refractive index of the lens is labeled as \textbackslash{}( n = 1.50 \textbackslash{}).
   - \textbackslash{}( R\_1 = \textbackslash{}infty \textbackslash{}) is noted, indicating one surface of the lens is flat.
   - \textbackslash{}( R\_2 \textbackslash{}) is denoted, possibly referring to the radius of curvature of the other surface.

3. **Light Rays**:
   - Two purple lines represent light rays, emanating from a small green arrow (the object).
   - The rays pass through the lens and converge at a point, creating an image.

4. **Text**: 
   - "Image and object distances not to scale" is noted, indicating that the diagram does not represent actual distances accurately.

5. **Arrows and Direction**:
   - The green arrow at the bottom suggests the direction of the object.
   - Light rays are depicted as moving from the object, through the lens, to the image.

This diagram represents the behavior of light as it passes through a lens to form an image at a different location.

\#\# Dataset metadata

Geometrical Optics; Physical Model Grounding Reasoning, Multi-Formula Reasoning

\paragraph{Recorded answer/context.}
C: C: 3.8mm

\paragraph{Figure/image path.}
/root/proposal\_for\_physic/science-mango/phyx\_mini\_run/phyx\_data/test\_image/57.png

\paragraph{Formalization target.}
create a compiling Lean file with sorry bodies at `PhyXMiniProblems/problem\_phyx\_mini\_0057.lean`. The Lean declarations must preserve the physical quantities, dimensions or dimensional roles, figure labels, governing-law hypotheses, and final relation expressed by this problem.
Use Mathlib/Physlib names found through LeanExplore where available. If a physics API is missing, introduce faithful local abstractions rather than scalar placeholder aliases.

\begin{theorem}[Physics formalization target]
\label{thm:physics:phyx_mini_0057:target}
\lean{PhyXMiniProblems.ProblemPhyXMini0057.problem_phyx_mini_0057}
\uses{def:physics:phyx-mini-0057:phyxminiproblems-problemphyxmini0057-lengthinmillimeters, def:physics:phyx-mini-0057:phyxminiproblems-problemphyxmini0057-microscopeobjectivesetup, def:physics:phyx-mini-0057:phyxminiproblems-problemphyxmini0057-matchesproblemandfigurereadouts, def:physics:phyx-mini-0057:phyxminiproblems-problemphyxmini0057-hasphysicalparameters, def:physics:phyx-mini-0057:phyxminiproblems-problemphyxmini0057-modelsambientair, def:physics:phyx-mini-0057:phyxminiproblems-problemphyxmini0057-satisfiesparaxialobjectivelaws, def:physics:phyx-mini-0057:phyxminiproblems-problemphyxmini0057-matchesanswertonearesttenth}
The assigned autoformalize agent should translate this physics problem into Lean declarations in the covered file, with theorem and lemma proof bodies written as `by sorry`.
\end{theorem}
\begin{proof}
This is an autoformalization task, not a proof task. Produce statements that can later be proved by the physics prover without weakening the physical model.
\end{proof}

% --- Archon declaration topology begin ---
\section{Declaration topology}
\begin{definition}[Length Magnitude]
  \label{def:physics:phyx-mini-0057:phyxminiproblems-problemphyxmini0057-lengthmagnitude}
  \lean{PhyXMiniProblems.ProblemPhyXMini0057.LengthMagnitude}
  A nonnegative physical quantity carrying the dimension of length
\end{definition}
\begin{proof}
  This is the declared model definition of Length Magnitude.
\end{proof}
\begin{definition}[length In Millimeters]
  \label{def:physics:phyx-mini-0057:phyxminiproblems-problemphyxmini0057-lengthinmillimeters}
  \lean{PhyXMiniProblems.ProblemPhyXMini0057.lengthInMillimeters}
  \uses{def:physics:phyx-mini-0057:phyxminiproblems-problemphyxmini0057-lengthmagnitude}
  The numerical readout, in millimeters, of a physical length magnitude
\end{definition}
\begin{proof}
  This is the declared model definition of length In Millimeters.
\end{proof}
\begin{definition}[Axial Side]
  \label{def:physics:phyx-mini-0057:phyxminiproblems-problemphyxmini0057-axialside}
  \lean{PhyXMiniProblems.ProblemPhyXMini0057.AxialSide}
  Axial sides relative to propagation from the specimen toward the image
\end{definition}
\begin{proof}
  This is the declared model definition of Axial Side.
\end{proof}
\begin{definition}[Lens Surface Label]
  \label{def:physics:phyx-mini-0057:phyxminiproblems-problemphyxmini0057-lenssurfacelabel}
  \lean{PhyXMiniProblems.ProblemPhyXMini0057.LensSurfaceLabel}
  The surface labels printed in the figure
\end{definition}
\begin{proof}
  This is the declared model definition of Lens Surface Label.
\end{proof}
\begin{definition}[Figure Distance Label]
  \label{def:physics:phyx-mini-0057:phyxminiproblems-problemphyxmini0057-figuredistancelabel}
  \lean{PhyXMiniProblems.ProblemPhyXMini0057.FigureDistanceLabel}
  The two distance-arrow labels printed in the figure
\end{definition}
\begin{proof}
  This is the declared model definition of Figure Distance Label.
\end{proof}
\begin{definition}[Optical Medium]
  \label{def:physics:phyx-mini-0057:phyxminiproblems-problemphyxmini0057-opticalmedium}
  \lean{PhyXMiniProblems.ProblemPhyXMini0057.OpticalMedium}
  The optical media relevant to the objective
\end{definition}
\begin{proof}
  This is the declared model definition of Optical Medium.
\end{proof}
\begin{definition}[Image Nature]
  \label{def:physics:phyx-mini-0057:phyxminiproblems-problemphyxmini0057-imagenature}
  \lean{PhyXMiniProblems.ProblemPhyXMini0057.ImageNature}
  Whether rays actually meet at the image or only appear to do so
\end{definition}
\begin{proof}
  This is the declared model definition of Image Nature.
\end{proof}
\begin{definition}[Propagation Direction]
  \label{def:physics:phyx-mini-0057:phyxminiproblems-problemphyxmini0057-propagationdirection}
  \lean{PhyXMiniProblems.ProblemPhyXMini0057.PropagationDirection}
  Direction in which the pictured rays propagate
\end{definition}
\begin{proof}
  This is the declared model definition of Propagation Direction.
\end{proof}
\begin{definition}[Lens Surface Profile]
  \label{def:physics:phyx-mini-0057:phyxminiproblems-problemphyxmini0057-lenssurfaceprofile}
  \lean{PhyXMiniProblems.ProblemPhyXMini0057.LensSurfaceProfile}
  \uses{def:physics:phyx-mini-0057:phyxminiproblems-problemphyxmini0057-lengthmagnitude, def:physics:phyx-mini-0057:phyxminiproblems-problemphyxmini0057-axialside}
  This inductive records the Lens Surface Profile component of the problem-specific physical model
\end{definition}
\begin{proof}
  This is the declared model definition of Lens Surface Profile.
\end{proof}
\begin{definition}[Planoconvex Objective]
  \label{def:physics:phyx-mini-0057:phyxminiproblems-problemphyxmini0057-planoconvexobjective}
  \lean{PhyXMiniProblems.ProblemPhyXMini0057.PlanoconvexObjective}
  \uses{def:physics:phyx-mini-0057:phyxminiproblems-problemphyxmini0057-lengthmagnitude, def:physics:phyx-mini-0057:phyxminiproblems-problemphyxmini0057-axialside, def:physics:phyx-mini-0057:phyxminiproblems-problemphyxmini0057-lenssurfacelabel, def:physics:phyx-mini-0057:phyxminiproblems-problemphyxmini0057-lenssurfaceprofile}
  The planoconvex objective together with its two labeled faces
\end{definition}
\begin{proof}
  This is the declared model definition of Planoconvex Objective.
\end{proof}
\begin{definition}[Objective Lens Figure]
  \label{def:physics:phyx-mini-0057:phyxminiproblems-problemphyxmini0057-objectivelensfigure}
  \lean{PhyXMiniProblems.ProblemPhyXMini0057.ObjectiveLensFigure}
  \uses{def:physics:phyx-mini-0057:phyxminiproblems-problemphyxmini0057-lengthmagnitude, def:physics:phyx-mini-0057:phyxminiproblems-problemphyxmini0057-figuredistancelabel}
  Quantitative arrows and qualitative ray information read from the figure
\end{definition}
\begin{proof}
  This is the declared model definition of Objective Lens Figure.
\end{proof}
\begin{definition}[Microscope Objective Setup]
  \label{def:physics:phyx-mini-0057:phyxminiproblems-problemphyxmini0057-microscopeobjectivesetup}
  \lean{PhyXMiniProblems.ProblemPhyXMini0057.MicroscopeObjectiveSetup}
  \uses{def:physics:phyx-mini-0057:phyxminiproblems-problemphyxmini0057-lengthmagnitude, def:physics:phyx-mini-0057:phyxminiproblems-problemphyxmini0057-axialside, def:physics:phyx-mini-0057:phyxminiproblems-problemphyxmini0057-opticalmedium, def:physics:phyx-mini-0057:phyxminiproblems-problemphyxmini0057-imagenature, def:physics:phyx-mini-0057:phyxminiproblems-problemphyxmini0057-propagationdirection, def:physics:phyx-mini-0057:phyxminiproblems-problemphyxmini0057-planoconvexobjective, def:physics:phyx-mini-0057:phyxminiproblems-problemphyxmini0057-objectivelensfigure}
  Physical quantities and roles in the thin-lens microscope setup
\end{definition}
\begin{proof}
  This is the declared model definition of Microscope Objective Setup.
\end{proof}
\begin{definition}[surface Curvature In Inverse Millimeters]
  \label{def:physics:phyx-mini-0057:phyxminiproblems-problemphyxmini0057-surfacecurvatureininversemillimeters}
  \lean{PhyXMiniProblems.ProblemPhyXMini0057.surfaceCurvatureInInverseMillimeters}
  \uses{def:physics:phyx-mini-0057:phyxminiproblems-problemphyxmini0057-lengthinmillimeters, def:physics:phyx-mini-0057:phyxminiproblems-problemphyxmini0057-lenssurfaceprofile}
  This def records the surface Curvature In Inverse Millimeters component of the problem-specific physical model
\end{definition}
\begin{proof}
  This is the declared model definition of surface Curvature In Inverse Millimeters.
\end{proof}
\begin{definition}[Matches Problem And Figure Readouts]
  \label{def:physics:phyx-mini-0057:phyxminiproblems-problemphyxmini0057-matchesproblemandfigurereadouts}
  \lean{PhyXMiniProblems.ProblemPhyXMini0057.MatchesProblemAndFigureReadouts}
  \uses{def:physics:phyx-mini-0057:phyxminiproblems-problemphyxmini0057-lengthinmillimeters, def:physics:phyx-mini-0057:phyxminiproblems-problemphyxmini0057-microscopeobjectivesetup}
  This def records the Matches Problem And Figure Readouts component of the problem-specific physical model
\end{definition}
\begin{proof}
  This is the declared model definition of Matches Problem And Figure Readouts.
\end{proof}
\begin{definition}[Has Physical Parameters]
  \label{def:physics:phyx-mini-0057:phyxminiproblems-problemphyxmini0057-hasphysicalparameters}
  \lean{PhyXMiniProblems.ProblemPhyXMini0057.HasPhysicalParameters}
  \uses{def:physics:phyx-mini-0057:phyxminiproblems-problemphyxmini0057-lengthinmillimeters, def:physics:phyx-mini-0057:phyxminiproblems-problemphyxmini0057-microscopeobjectivesetup}
  Positivity and refractive-index branch conditions for the physical setup
\end{definition}
\begin{proof}
  This is the declared model definition of Has Physical Parameters.
\end{proof}
\begin{definition}[Models Ambient Air]
  \label{def:physics:phyx-mini-0057:phyxminiproblems-problemphyxmini0057-modelsambientair}
  \lean{PhyXMiniProblems.ProblemPhyXMini0057.ModelsAmbientAir}
  \uses{def:physics:phyx-mini-0057:phyxminiproblems-problemphyxmini0057-microscopeobjectivesetup}
  The surrounding optical medium is ordinary air, of refractive index one
\end{definition}
\begin{proof}
  This is the declared model definition of Models Ambient Air.
\end{proof}
\begin{definition}[Satisfies Paraxial Objective Laws]
  \label{def:physics:phyx-mini-0057:phyxminiproblems-problemphyxmini0057-satisfiesparaxialobjectivelaws}
  \lean{PhyXMiniProblems.ProblemPhyXMini0057.SatisfiesParaxialObjectiveLaws}
  \uses{def:physics:phyx-mini-0057:phyxminiproblems-problemphyxmini0057-lengthinmillimeters, def:physics:phyx-mini-0057:phyxminiproblems-problemphyxmini0057-microscopeobjectivesetup, def:physics:phyx-mini-0057:phyxminiproblems-problemphyxmini0057-surfacecurvatureininversemillimeters}
  This structure records the Satisfies Paraxial Objective Laws component of the problem-specific physical model
\end{definition}
\begin{proof}
  This is the declared model definition of Satisfies Paraxial Objective Laws.
\end{proof}
\begin{definition}[Answer Choice]
  \label{def:physics:phyx-mini-0057:phyxminiproblems-problemphyxmini0057-answerchoice}
  \lean{PhyXMiniProblems.ProblemPhyXMini0057.AnswerChoice}
  Labels of the four displayed multiple-choice answers
\end{definition}
\begin{proof}
  This is the declared model definition of Answer Choice.
\end{proof}
\begin{definition}[answer Radius In Millimeters]
  \label{def:physics:phyx-mini-0057:phyxminiproblems-problemphyxmini0057-answerradiusinmillimeters}
  \lean{PhyXMiniProblems.ProblemPhyXMini0057.answerRadiusInMillimeters}
  \uses{def:physics:phyx-mini-0057:phyxminiproblems-problemphyxmini0057-answerchoice}
  The radius readout printed beside each answer label, in millimeters
\end{definition}
\begin{proof}
  This is the declared model definition of answer Radius In Millimeters.
\end{proof}
\begin{definition}[Matches Answer To Nearest Tenth]
  \label{def:physics:phyx-mini-0057:phyxminiproblems-problemphyxmini0057-matchesanswertonearesttenth}
  \lean{PhyXMiniProblems.ProblemPhyXMini0057.MatchesAnswerToNearestTenth}
  \uses{def:physics:phyx-mini-0057:phyxminiproblems-problemphyxmini0057-answerchoice, def:physics:phyx-mini-0057:phyxminiproblems-problemphyxmini0057-answerradiusinmillimeters}
  Agreement of an exact radius with a displayed nearest-tenth readout
\end{definition}
\begin{proof}
  This is the declared model definition of Matches Answer To Nearest Tenth.
\end{proof}
\begin{lemma}[surface Curvature Readouts]
  \label{lem:physics:phyx-mini-0057:phyxminiproblems-problemphyxmini0057-surfacecurvaturereadouts}
  \lean{PhyXMiniProblems.ProblemPhyXMini0057.surfaceCurvatureReadouts}
  \uses{def:physics:phyx-mini-0057:phyxminiproblems-problemphyxmini0057-lengthinmillimeters, def:physics:phyx-mini-0057:phyxminiproblems-problemphyxmini0057-microscopeobjectivesetup, def:physics:phyx-mini-0057:phyxminiproblems-problemphyxmini0057-surfacecurvatureininversemillimeters, def:physics:phyx-mini-0057:phyxminiproblems-problemphyxmini0057-matchesproblemandfigurereadouts}
  This lemma records the surface Curvature Readouts component of the problem-specific physical model
\end{lemma}
\begin{proof}
  The conclusion follows from the governing relations and figure readouts named in the statement of surface Curvature Readouts.
\end{proof}
\begin{lemma}[reciprocal Focal Length eq twenty one over one sixty]
  \label{lem:physics:phyx-mini-0057:phyxminiproblems-problemphyxmini0057-reciprocalfocallength-eq-twenty-one-over-one-sixty}
  \lean{PhyXMiniProblems.ProblemPhyXMini0057.reciprocalFocalLength_eq_twenty_one_over_one_sixty}
  \uses{def:physics:phyx-mini-0057:phyxminiproblems-problemphyxmini0057-lengthinmillimeters, def:physics:phyx-mini-0057:phyxminiproblems-problemphyxmini0057-microscopeobjectivesetup, def:physics:phyx-mini-0057:phyxminiproblems-problemphyxmini0057-matchesproblemandfigurereadouts, def:physics:phyx-mini-0057:phyxminiproblems-problemphyxmini0057-satisfiesparaxialobjectivelaws}
  The imaging equation and the two distance readouts give 1/f = 21/160
\end{lemma}
\begin{proof}
  The conclusion follows from the governing relations and figure readouts named in the statement of reciprocal Focal Length eq twenty one over one sixty.
\end{proof}
\begin{lemma}[curved Surface Radius eq eighty over twenty one]
  \label{lem:physics:phyx-mini-0057:phyxminiproblems-problemphyxmini0057-curvedsurfaceradius-eq-eighty-over-twenty-one}
  \lean{PhyXMiniProblems.ProblemPhyXMini0057.curvedSurfaceRadius_eq_eighty_over_twenty_one}
  \uses{def:physics:phyx-mini-0057:phyxminiproblems-problemphyxmini0057-lengthinmillimeters, def:physics:phyx-mini-0057:phyxminiproblems-problemphyxmini0057-microscopeobjectivesetup, def:physics:phyx-mini-0057:phyxminiproblems-problemphyxmini0057-matchesproblemandfigurereadouts, def:physics:phyx-mini-0057:phyxminiproblems-problemphyxmini0057-hasphysicalparameters, def:physics:phyx-mini-0057:phyxminiproblems-problemphyxmini0057-modelsambientair, def:physics:phyx-mini-0057:phyxminiproblems-problemphyxmini0057-satisfiesparaxialobjectivelaws}
  This lemma records the curved Surface Radius eq eighty over twenty one component of the problem-specific physical model
\end{lemma}
\begin{proof}
  The conclusion follows from the governing relations and figure readouts named in the statement of curved Surface Radius eq eighty over twenty one.
\end{proof}
% --- Archon declaration topology end ---
% --- Archon physics formalization source end ---
